% Overleaf 업로드용 — 컴파일러: XeLaTeX 권장 (Menu > Compiler > XeLaTeX).
% pdfLaTeX로도 kotex가 동작하지만 한글 폰트 문제 시 XeLaTeX로 바꾸세요.
% 출처: docs/KJH/screen-design.md (2026-07-14 기준). 확정 사양은 레포 문서가 기준.
\documentclass[11pt]{article}
\usepackage{kotex}
\usepackage[a4paper,margin=2.3cm]{geometry}
\usepackage{enumitem}
\usepackage{xcolor}
\usepackage{titlesec}
\usepackage{hyperref}
\hypersetup{colorlinks=true, linkcolor=black, urlcolor=blue}
\setlist{itemsep=2pt, topsep=3pt}
\titleformat{\section}{\Large\bfseries}{\thesection}{0.6em}{}
\titleformat{\subsection}{\large\bfseries}{\thesubsection}{0.6em}{}

\definecolor{warnbg}{HTML}{FFF3CD}
\definecolor{warnbd}{HTML}{E0A800}
\newcommand{\tbd}{\textcolor{red}{\textbf{[미확정]}}}

\newcommand{\warnbox}[1]{%
  \begin{center}\fcolorbox{warnbd}{warnbg}{%
  \parbox{0.95\linewidth}{\vspace{2pt}#1\vspace{2pt}}}\end{center}}

\title{\textbf{게임 프론트엔드 · 규칙 사양서}\\[4pt]\large 멀티플레이 AI 릴레이 맵메이커 (초안)}
\author{KJH → 팀원 공유용}
\date{2026-07-14}

\begin{document}
\maketitle

\warnbox{%
\textbf{먼저 읽기 — 문서 상태}
\begin{itemize}[leftmargin=1.4em]
  \item \textbf{API 스펙은 전부 미정.} 엔드포인트·요청/응답 형태가 확정되지 않았습니다. 화면을 먼저 만들되 \textbf{서버 연동부는 인터페이스로 추상화}하고 경로·페이로드를 \textbf{하드코딩하지 마세요.}
  \item \textbf{타일당 px(tileSize)는 바뀔 가능성이 높습니다.} 모든 길이·간격·크기는 \textbf{tileSize 기준 상수/비율}로. px 하드코딩 금지.
  \item \tbd{} 태그 = 아직 확정 안 된 항목. 가정을 두되 바뀔 수 있음을 전제.
  \item \textbf{디자인 참고:} \emph{Super Mario Maker 2} 유튜브 플레이 영상을 보며 UI·연출 톤 참고. 단 아래 ``UI 테마/모션'' 원칙(성숙하게 렌더링) 기준.
\end{itemize}}

\section{담당 분담}
\begin{itemize}
  \item \textbf{팀원(프론트 전반):} 로그인 · 메인 · 아바타 생성 · 에셋 스튜디오 · 창고 · 설정 · 로비 · 방 대기실, 그리고 \textbf{레이스 화면(E) — 가능하면 시도}, \textbf{결과 화면(F) — 요청}.
  \item \textbf{KJH 전담:} \textbf{맵 에디터(S4)} 는 KJH가 직접 구현. 이 문서엔 이해용 개요·연동점만. (팀원 구현 대상 아님)
  \item AI 백엔드(LLM · 생성 파이프라인 · 큐)도 전부 KJH.
\end{itemize}

\section{전체 플로우}
로그인 $\rightarrow$ 메인 $\rightarrow$ (아바타 제작 / 에셋 스튜디오) $\rightarrow$ 로비 $\rightarrow$ 방 $\rightarrow$ 제작(3분) $\rightarrow$ 검증(테스트) $\rightarrow$ 레이스 $\rightarrow$ 결과.

\section{게임 규칙 (레이스)}
\warnbox{아래 수치·규칙은 기획 확정분이나 \textbf{밸런싱으로 조정될 수 있음.} 모든 길이/시간은 상수로.}

\subsection{방 (로비 → 방 대기실)}
\begin{itemize}
  \item \textbf{방 생성:} 공개방 / 비밀방. 비밀방은 \textbf{비밀번호 설정} 필요. 그 외 설정 없음.
  \item \textbf{입장:} 공개방 목록에서 입장, 또는 비밀방 비번 입력.
  \item \textbf{난입:} 맵 제작 페이즈 동안 언제든 난입 가능. 단 \textbf{제작시간 1분 미만} 남은 시점에 난입하면 그 플레이어는 \textbf{맵 제작 불가}(대기/관전).
\end{itemize}

\subsection{제작 페이즈 (3분)}
\begin{itemize}
  \item 각 플레이어가 \textbf{자기 라인 1개}를 제작.
  \item \textbf{사용 에셋 제한:}
  \begin{itemize}
    \item 내가 만든 에셋: 카테고리 무관 \textbf{최대 3개.}
    \item DB(공용) 에셋: 카테고리별 정해진 개수만큼 \textbf{랜덤 제공.}
    \item 배경 에셋: 예외 — \textbf{제한 없음, 자유 사용.}
  \end{itemize}
  \item \textbf{15초 프리뷰}(3분 시작 전): 사용 가능 에셋이 화면에 뜸. \textbf{다른 유저가 만든 것도} 봄. 카테고리·\textbf{공용/개인 구분} 표시, hover 시 상세.
  \item 15초 후 \textbf{맵 에디터} 진입.
\end{itemize}

\subsection{라인 · 깃발}
\begin{itemize}
  \item \textbf{라인 최대 길이 고정}(상수). 실제 길이는 시작/끝 깃발 위치로 \textbf{유동 결정.}
  \item 깃발 = \textbf{마리오 체크포인트식.} 에디터 열면 \textbf{미리 대충 배치}돼 있음.
  \item 깃발 아래 \textbf{3$\times$1 바닥 = 깃발과 한 세트.}
  \item \textbf{끝 깃발은 시작 깃발보다 왼쪽 불가.} 둘 다 가로·세로 자유 이동.
  \item 기본 \textbf{완전 빈 공간}, 두 깃발 아래 3타일만 기본 땅.
  \item \textbf{테스트:} 라인을 \textbf{1회 이상 테스트 완료 못 하면 맵에 안 쓰임.} \textbf{가장 최근 완료된 테스트 라인}이 사용됨.
\end{itemize}

\subsection{병합 · 골}
\begin{itemize}
  \item 제작 종료 후 \textbf{테스트 완료 라인끼리 랜덤 순서로 연결.}
  \item \textbf{맨 첫 라인의 진짜 시작 / 맨 끝 라인의 진짜 골} 깃발만 \textbf{깃대를 길게}(구분).
  \item 마지막 라인 끝에 \textbf{골 지점.}
\end{itemize}

\subsection{레이스}
\begin{itemize}
  \item \textbf{게임 시간 = 라인 수 $\times$ 40초.}
  \item 먼저 \textbf{골대를 잡은 사람이 1등.} 1등 도달 시 \textbf{10초 카운트다운} → 그 안에 들어온 사람은 \textbf{순서대로 순위}, 나머지는 \textbf{리타이어.}
  \item 게임 시간 종료까지 1등이 없으면 → \textbf{30초 ``라스트댄스''} → 30초 뒤 \textbf{가장 앞선 순서로 1/2/3등만} 순위.
  \begin{itemize}
    \item 라스트댄스 동안 \textbf{1/2/3위 머리 위에 금/은/동 긴 세로선}(직관적 판단).
  \end{itemize}
  \item \textbf{라인 파괴 스윕:} 게임 시간 동안 \textbf{첫 라인부터 순차로 50초마다} 그 라인의 모든 에셋이 파괴되고 시작·골 지점 땅이 쭉 이어짐(\textbf{라인 제작자가 유리해지는 것 방지}).
\end{itemize}

\subsection{라인 저장 (백엔드)}
라인은 \textbf{테스트 성공 시 영속 저장}(MapLine 엔티티, \tbd{} 스키마 설계 중). 나중에 ``라인 카테고리''로 남의 라인 재사용 대비.

\section{UI 테마 (요약)}
\begin{itemize}
  \item 원칙: SMM2의 색·격자 모티프를 \textbf{콘솔 런처 그라데이션 질감으로 성숙하게} 렌더링. 유아스러운 표현(두꺼운 검정 아웃라인, 전면 도트 폰트, 원색 단색 면) 금지.
  \item 색: 건설 노랑 \texttt{\#F6BE00}(주역) / 스카이 블루 \texttt{\#4A9DE0}(배경 그라데이션) / 마리오 레드 \texttt{\#E52521}(포인트) / 그린·브라운(보조). \emph{hex는 조정 가능한 기준값.}
  \item 질감: 은은한 수직 그라데이션 + 부드러운 광택, 살짝만 라운드. 색면 위 \textbf{아주 투명한 흰 타일 격자 오버레이}(rgba(255,255,255,0.08) 안팎)를 일관 적용.
  \item 폰트: 타이틀·타이머 = 도트/블록형(Galmuri) / 본문·버튼 = Pretendard.
\end{itemize}

\section{UI 모션 · 마이크로인터랙션}
\warnbox{원칙: \textbf{선형(linear) 이동 금지} — 모든 등장·반응에 스프링/이징. ``생동감 있되 유아스럽지 않게.''}
\begin{itemize}
  \item \textbf{등장:} 버튼·카드·패널이 아래→위로 \textbf{스프링 오버슈트}(살짝 튀어올랐다 정착). 동시 등장은 \textbf{스태거}(0.03\textasciitilde0.06s 간격).
  \item \textbf{호버:} 미세 확대(scale 1.02\textasciitilde1.05) + 은은한 광택/그림자 상승.
  \item \textbf{클릭:} 눌림(scale $\approx$0.94) → \textbf{스프링 복귀로 ``띠용''}(오버슈트 바운스) + 색·그림자 짧은 펄스.
  \item \textbf{화면 전환:} 페이드 + 슬라이드를 \textbf{이징(cubic-bezier, 선형 아님)}으로.
  \item \textbf{상태 모션:} 생성 중 프로그레스·타이머 등 진행감 있는 미세 모션.
\end{itemize}
\textbf{구현 원칙:}
\begin{itemize}
  \item \textbf{모션 토큰화:} 즉흥값 대신 공통 토큰 세트(예: \texttt{spring.pop=\{stiffness:500,damping:28\}}, \texttt{ease.enter=cubic-bezier(0.22,1,0.36,1)}).
  \item \textbf{접근성:} \texttt{prefers-reduced-motion} 존중(모션 최소화 대체).
  \item \textbf{라이브러리:} React+Vite → \textbf{framer-motion} 또는 react-spring 권장. 선택은 팀원 재량 \tbd.
\end{itemize}

\section{화면 사양}

\subsection{S1. 로그인 (확정)}
닉네임 1필드(1\textasciitilde12자) + [시작하기]. \texttt{POST /session}(\tbd) → token 발급, localStorage 저장. 재방문 시 token 있으면 메인 직행. 닉네임 중복 허용(식별=token).

\subsection{S2. 메인 (확정)}
버튼: [게임하기]→로비 / [에셋 만들기]→스튜디오 / [내 창고] / 우상단 ⚙️→설정. 좌측 아바타 패널(미보유=졸라맨, 생성 중 상태 표시). 아바타 강제 안 함. \textbf{아바타 = category=avatar 에셋.}

\subsection{S2b. 내 창고 (확정)}
2단 탭 [아바타]/[컴포넌트 에셋](내부 카테고리 소필터). 썸네일 그리드. 생성 상태 오버레이(queued/generating/ready/failed). 클릭 → 에셋 상세/검수 팝업(애니메이션 미리보기 + 액션별 재생성). 상태 갱신은 소켓 푸시 우선 \tbd.

\subsection{S2c. 설정 모달 (확정)}
소리(볼륨/음소거, localStorage) · 닉네임 변경(DB) · 기기 연동(단축 코드 방식, 5분 유효).

\subsection{A. 아바타 생성 (확정)}
좌 도구 패널 + 중앙 캔버스(1$\times$2 타일, 256$\times$512, 투명) + 하단 설명/생성. 도구: 펜·지우개·굵기·스포이드·전체이동·실행취소·전체지우기. \textbf{이미지 업로드 없음}(손그림만). [생성하기] → \texttt{POST /assets}(category=avatar, \tbd) → 즉시 메인 복귀, 진행은 아바타 패널이 표시.

\subsection{B. 에셋 스튜디오 (확정)}
그림판 도구·캔버스는 A와 공통(단 캔버스 비율 = 선택 크기 n$\times$m). 좌 도구+팔레트+[에셋 불러오기] / 우 카테고리+크기+속성 / 중앙 캔버스. 애니메이션은 \textbf{카테고리+옵션이 자동 결정·생성}(유저 개별 토글 안 함). 이름 필수.

\subsection{S3. 로비 (초안, 팀원)}
공개방 카드 목록 + [방 생성]/[공개방 입장]/비밀방 비번. 카드: 방장 닉네임 · 현/총 인원 · 플레이시간(진행 경과) · 공개/비밀. 진행 중(레이스) 방 입장 불가. \tbd{} 정렬·갱신 방식·API.

\subsection{C. 방 대기실 (초안, 팀원)}
공개/비밀 선택·비번 설정. 참가자 목록 + 방장 [시작]. 제작 페이즈 중 난입 가능(1분 미만 시 제작 불가). \tbd{} 최소 인원·API.

\subsection{S4. 맵 에디터 (개요, KJH 전담)}
\textbf{팀원 구현 대상 아님 — 흐름 이해 + 연동점(에셋 목록·라인 저장 API, 둘 다 \tbd)만.}
레이아웃: 좌 에셋 창고(내가만든/남이만든 · 카테고리 필터 · 검색 · 결과 리스트) / 상단 즐겨찾기 실사용 슬롯(내 3개 + 랜덤 공용) / 중앙 캔버스(\textbf{1타일 스냅}) / 우 툴바 \tbd / 우하단 시간추가·단축·남은시간·[테스트하기]. 깃발·라인 규칙은 ``게임 규칙'' 참조. 이동/스위치/경사 배치는 추가 데이터 필요 \tbd.

\subsection{D. 검증(테스트) (개요, KJH/공용)}
[테스트하기]로 자기 라인 단독 플레이 테스트. \textbf{완주 성공 시}만 ``테스트 완료''로 저장(최근 성공본). 미완료 라인은 맵에 안 쓰임. \tbd{} 진입/복귀 UI.

\subsection{E. 레이스 (초안, 팀원 — 가능하면)}
병합 맵 동시 레이스(좌→우). HUD: 남은시간(라인수$\times$40s)·실시간 순위·1등 후 10초 카운트다운·라스트댄스(금/은/동 세로선)·50초 라인 파괴 스윕 연출. \textbf{물리·네트워크는 서버 권위(Colyseus)+공유 물리, 렌더는 Phaser} → 인게임 로직은 KJH쪽과 맞물림(팀원은 렌더/HUD 중심). \tbd{} HUD 상세·상태 동기 계약·API.

\subsection{F. 결과 (초안, 팀원 — 요청)}
최종 순위표(1/2/3등 금/은/동 + 리타이어). 플레이어별 기록 \tbd. [다시하기]/[로비로]. \tbd{} 표시 데이터·API.

\section{데이터 · 연동 (개요)}
\warnbox{\textbf{모든 API는 아직 \tbd{}.} 아래는 개념 모델일 뿐, 연동부는 인터페이스로 감싸 두세요.}
\begin{itemize}
  \item \textbf{Asset}: 유저/시스템 제작 에셋(아바타 포함). category, attrs(옵션), 상태(생성 큐).
  \item \textbf{AssetSprite}: 에셋의 액션별 스프라이트 시트(idle/walk/onair 등) — AI 생성 파이프라인 산출.
  \item \textbf{MapLine} \tbd: 제작·테스트 완료된 라인(제작자·배치·깃발·길이·테스트완료여부·공개여부).
  \item \textbf{동기화}: 서버 권위 + Colyseus(실시간). 에셋 생성 상태는 소켓 푸시(폴링 대안) \tbd.
\end{itemize}

\end{document}
